\section{Autoencoder-Based Representation Learning in Alzheimer's Disease}
Autoencoder-based representation learning has been used to identify patterns in neuroimaging data related to Alzheimer’s disease (AD) and mild cognitive impairment (MCI). 
Early work by Suk and Shen used autoencoders to learn features from neuroimaging and biomarker data for AD/MCI classification~\cite{SukShen2013}. 
Their results showed that learned representations can contain disease-related information useful for distinguishing between clinical groups, providing an alternative to manually defined imaging features.

Autoencoders have also been applied to resting-state fMRI (rs-fMRI) to represent functional brain activity. 
Suk et al.\ used a deep autoencoder to map regional BOLD activity at individual timepoints into a lower-dimensional representation~\cite{SukStateSpace}. 
Applying the encoder across the rs-fMRI time series produced a sequence of latent representations, which was then used to model functional dynamics and distinguish individuals with MCI from healthy controls. 
Other approaches have instead learned representations from functional connectivity derived from rs-fMRI.\@ 
For example, autoencoder-based feature learning has been applied to functional brain networks for the early identification of MCI~\cite{EarlyDiagnosis2018}. 
These approaches show that disease-related information can be learned both directly from regional brain activity and from the functional relationships between brain regions.

More recent work has extended this approach towards probabilistic latent representations. 
Variational methods have been used to represent rs-fMRI functional connectivity through latent distributions for the identification of MCI~\cite{FCLatentDistributions2023}.
This allows functional connectivity patterns to be represented in a structured latent space rather than only through deterministic features. 
Overall, previous studies show that autoencoders can learn AD- and MCI-related representations from functional neuroimaging data. 
However, these representations have mainly been used as features for identifying or classifying clinical groups, with less focus on how disease-related differences are represented within the latent space and reconstructed back into functional brain activity.